Dado el siguiente problema de programación lineal:

\begin{equation}
\begin{split}
\mbox {max. }z=2x_1+1x_2+3x_3\\
s.a.:\\
4x_1+2x_2+1x_3 \leq 60\\
2x_1+6x_2 -1x_3 \leq 20\\
1x_1+2x_3 \leq 10\\
x_1,\,x_2,\,x_3 \geq 0
\notag
\end{split}
\end{equation}

Se pide
\begin{enumerate}
    \item \textbf{Identificar la base correspondiente a la solución  en la que $x_1=0$, $x_2=25/6$, $x_3=4$.}

    El problema en formato estándar es:
    
    \begin{equation}
    \begin{split}
    \mbox {max. }z=2x_1+1x_2+3x_3\\
    s.a.:\\
    4x_1+2x_2+1x_3 + h_1 = 60\\
    2x_1+6x_2 -1x_3 + h_2 = 20\\
    1x_1+2x_3 + h_3 = 10\\
    x_1,\,x_2,\,x_3,\,h_1,\,h_2,\,h_3 \geq 0
    \notag
    \end{split}
    \end{equation}

    Despejando de las restricciones se obtiene que $h_1=140/3$, $h_2=0$ y $h_3=0$.

    Hay tres variables no nulas que son las variables básicas, $x_2$, $x_3$ y $h_1$, y las correspondientes columnas de la matríz A son las siguientes, en el mismo orden:

    \[
    \begin{pmatrix}
       2 &  1 &  1\\
       6 & -1 &  0\\
       0 &  2 &  0
    \end{pmatrix} 
    \]
    
    \item \textbf{Identificar la base correspondiente a la solución  en la que $x_1=10$, $x_2=0$, $x_3=0$.}

    Despejando de las restricciones se obtiene que $h_1=20$, $h_2=0$ y $h_3=0$.

    Hay dos variables no nulas, $x_1$ y $h_1$, que son dos variables básicas y las columnas en la matriz A correspondientes a estas dos variables serán dos de las columnas de la base.

    \[
            \begin{pmatrix}
               4 & 1 &  * \\
               2 & 0 &  *\\
               1 & 0 &  * 
            \end{pmatrix} 
            \]
    
    Se puede seleccionar otra variable nula de entre las restantes y las columnas de las tres variables componen cada una de las bases que conducen a la misma solución degenerada.
    
    \begin{itemize}
        \item $x_2$ variable básica. La tercera es la correspondiente a $x_2$ en A:
            \[
            \begin{pmatrix}
               4 &  1 &  2\\
               2 &  0 &  6 \\
               1 &  0 &  1 
            \end{pmatrix} 
            \]
        \item $h_2$ variable básica. La tercera es la correspondiente a $h_2$ en A:
        \[
            \begin{pmatrix}
               4 &  1 &  0\\
               2 &  0 &  1\\
               1 &  0 &  0
            \end{pmatrix} 
            \]
        
        \item $h_3$ variable básica. La tercera es la correspondiente a $h_3$ en A:

        \[
            \begin{pmatrix}
               4 &  1 &  0\\
               2 &  0 &  0\\
               1 &  0 &  1
            \end{pmatrix} 
            \]
    \end{itemize}
    
\end{enumerate}